\chapter{The FRI Protocol}
\label{ch:fri}
 
This chapter presents the Fast Reed-Solomon Interactive Oracle Proof of
Proximity, the protocol at the core of the system implemented in
this thesis. We begin by stating the problem FRI solves, then describe
the algebraic structure of the evaluation domain, and finally the
folding operation that drives the protocol.

\section{Overview and Problem Statement}
\label{sec:fri_overview}

 
Recall from \cref{def:rs_code} that, for a subset $S \subseteq \F$ and
a rate parameter $\rho \in (0,1)$, the Reed-Solomon code
$\RS[\F, S, \rho]$ consists of the evaluations on $S$ of polynomials of
degree less than $\rho|S|$. The problem FRI addresses is the following.
 
\begin{definition}[Rate Proximity Testing (RPT) {\cite[Section~2.1]{BBHR18}}]
\label{def:rpt}
A verifier is given oracle access to a function $f : S \to \F$, together
with auxiliary information supplied by a proer. The verifier must
distinguish, with high probability and using only a small number of
queries, between the case that $f \in \RS[\F, S, \rho]$ and the case
that $f$ is far from every member of $\RS[\F, S, \rho]$ in relative
Hamming distance (\cref{def:hamming}).
\end{definition}
 
RPT is a special case of the low-degree testing problem, and
Ben-Sasson et al.\ identify it as a major bottleneck for transparent
proof systems~\cite[Section~2.1]{BBHR18}.
 
A naive solution is instructive. The verifier could ask the prover to
open $f$ at $\rho|S|$ positions, interpolate the unique polynomial of
degree less than $\rho|S|$ through them, and check that it agrees with
$f$ elsewhere. This fails on two counts: the number of openings is
proportional to the degree bound, which can reach $2^{20}$ or more in
practical instances, and every opening carries a Merkle authentication
path of $\log|S|$ hash values (\cref{sec:merkle_trees}), so the proof
size becomes prohibitive.
 
FRI takes a different route. Instead of testing the degree directly,
it \emph{reduces} it: the prover repeatedly halves the degree of the
committed polynomial using randomness supplied by the verifier, until
the degree is small enough to send the remaining polynomial directly.
Each reduction step is committed, and the verifier spot-checks the
consistency of the steps at a small number of random positions. If the
original function is far from the code, then with high probability at
least one reduction is inconsistent, and the verifier detects it.
 
This yields a protocol with prover arithmetic complexity strictly
linear in $|S|$ and verifier arithmetic complexity strictly
logarithmic~\cite[Section~2.1]{BBHR18}. In the terminology of
\cref{sec:iop}, FRI is an interactive oracle proof of proximity (IOPP)
for Reed-Solomon codes: the prover's messages are committed
codewords which the verifier queries at a bounded number of positions,
rather than reading in full.
 
\begin{remark}
\label{rem:proximity_not_membership}
FRI establishes \emph{proximity}, not membership. A successful
verification does not certify that the committed vector is exactly a
Reed-Solomon codeword, only that it is close to one in relative
Hamming distance. This distinction is central to the soundness
analysis of the FRI protocol.
\end{remark}  % citare soundess proof
 

 
% ==================================================================
\section{The Evaluation Domain}
\label{sec:fri_domain}
% ==================================================================
 
The folding step of FRI relies on specific algebraic structure in the
evaluation domain. We state the requirements first, then give the
co(\cref{sec:fri_folding})nstruction used in our implementation.
 
 
\subsection{Requirements}
\label{subsec:domain_requirements}
 
The reduction from a polynomial $f$ to a polynomial of half the degree
is carried out by pairing the evaluations of $f$ at $x$ and at $-x$. %(\cref{sec:fri_folding})
For this to be possible at every position,
the domain must be closed under negation. Moreover, the folded
polynomial is committed not on $D$ but on the squared domain
$D^{(2)} := \{x^2 : x \in D\}$, and since the folding step is applied
iteratively, $D^{(2)}$ must be closed under negation as well, and so
on for every subsequent layer~\cite[Section~3.9.1]{ethSTARK}.
 
We therefore require a chain of domains
%
\[
D = D_0, \quad D_{i+1} = D_i^{(2)}, \quad i = 0, 1, \ldots
\]
%
such that every $D_i$ is closed under negation and
$|D_{i+1}| = |D_i|/2$. These requirements are met by taking $D$ to be
a coset of a multiplicative subgroup of order
$2^k$~\cite[Section~3.9.1]{ethSTARK}.
 
A coset is used rather than the subgroup itself so that the evaluation
domain is disjoint from the trace domain on which the polynomial being
committed is defined~\cite[Section~3.4]{ethSTARK}. Disjointness is
required in the full STARK protocol, where constraint polynomials are
expressed as rational functions whose denominators vanish on the trace
domain; evaluating them on a disjoint domain keeps those denominators
nonzero.
 
 
\subsection{Construction}
\label{subsec:domain_construction}

We build the domain from the roots of unity of \cref{subsec:roots_of_unity}.
Let $n$ be the number of coefficients of the polynomial to be committed,
$b$ the blowup factor, and $N = b \cdot n = 2^k$ the size of the
evaluation domain. Two elements are fixed:
%
\begin{equation}
\label{eq:domain_generators}
\omega = \gamma^{(p-1)/2^{k}}, \qquad
g = \gamma^{(p-1)/2^{k+1}},
\end{equation}
%
where $\gamma$ is a fixed multiplicative generator of $\Fp^*$. Here
$\omega$ is a primitive $2^k$-th root of unity: its powers
$\langle\omega\rangle = \{1, \omega, \ldots, \omega^{N-1}\}$ form a
cyclic group of $N$ equally spaced points. The element $g$ is a
primitive $2^{k+1}$-th root of unity, one ``half-step'' finer, and it
plays the role of an \emph{offset} that rotates the whole group. The
two are linked by the invariant
%
\begin{equation}
\label{eq:g_squared}
g^2 = \omega,
\end{equation}
%
which, as we show below, is exactly what makes the domain reproduce
itself under squaring.

The evaluation domain is the coset (\cref{def:coset})
%
\begin{equation}
\label{eq:domain}
D = g \cdot \langle \omega \rangle
  = \{\, g\omega^j : j = 0, \ldots, N-1 \,\},
\end{equation}
%
i.e.\ the group $\langle\omega\rangle$ shifted by the offset $g$. The
codeword committed by the prover is $\bigl(f(g\omega^j)\bigr)_{j=0}^{N-1}$,
computed by scaling the $i$-th coefficient of $f$ by $g^i$ and applying
the NTT (\cref{subsec:ntt,subsec:lde}). In our implementation $\omega$
and $g$ are returned by \texttt{two\_adic\_generator(k)} and
\texttt{two\_adic\_generator(k+1)} respectively, and the scaled NTT is
performed by \texttt{coset\_lde}.

Three properties follow, and together they justify the folding step.
 
\begin{proposition}
\label{prop:domain}
With $D$ as in \cref{eq:domain} and $N = 2^k$, $k \geq 1$:
\begin{enumerate}
  \item $\omega^{N/2} = -1$;
  \item $D$ is closed under negation: $-\,g\omega^{j} = g\omega^{j+N/2}$
    for $j < N/2$;
  \item $D^{(2)} = \omega \cdot \langle\omega^2\rangle$, a coset of the
    subgroup of order $N/2$; in particular $|D^{(2)}| = N/2$.
\end{enumerate}
\end{proposition}

\begin{proof}
\emph{(1)} The element $\omega^{N/2}$ squares to
$\omega^N = 1$, so it is a square root of unity. It cannot equal $1$,
since $\omega$ has order exactly $N$; as $\Fp$ is a field, the only
square roots of unity are $\pm 1$, hence $\omega^{N/2} = -1$.

\emph{(2)} Using (1),
$g\omega^{j+N/2} = g\omega^{j}\cdot\omega^{N/2} = -\,g\omega^{j}$, and
the index $j+N/2$ lies in $\{0,\dots,N-1\}$, so the negated element is
again in $D$.

\emph{(3)} Squaring an element of $D$ and applying \cref{eq:g_squared}
gives $(g\omega^{j})^2 = g^2\omega^{2j} = \omega^{2j+1}$, so
$D^{(2)} = \{\omega^{2j+1} : j = 0,\dots,N-1\}$, the subgroup
$\langle\omega^2\rangle$ shifted by $\omega$. By (2) the map
$x \mapsto x^2$ identifies $x$ with $-x$, so it is two-to-one on $D$
and $|D^{(2)}| = N/2$.
\end{proof}

Property~(3) is what makes the construction \emph{self-reproducing}:
the squared domain is again a coset of a two-power subgroup, of the
same shape as $D$ but half the size, now described by the pair
$(\omega, \omega^2)$ in place of $(g, \omega)$. A single FRI round
therefore updates the domain simply by squaring both parameters,
%
\begin{equation}
\label{eq:domain_update}
g \leftarrow g^2, \qquad \omega \leftarrow \omega^2,
\end{equation}
%
which in our implementation is the update
\texttt{offset = offset*offset}, \texttt{omega = omega*omega} applied
at the end of each layer. The requirements of
\cref{subsec:domain_requirements} are preserved at every level, as long
as the subgroup order remains even.

\begin{example}
\label{ex:domain_f17}
Take $\F_{17}$, whose multiplicative group has order $16 = 2^4$, with
generator $\gamma = 3$. For $N = 8$ we get $\omega = 3^2 = 9$ and
$g = 3^1 = 3$, so $g^2 = 9 = \omega$ as required by
\cref{eq:g_squared}. The evaluation domain is
%
\[
D = \{3, 10, 5, 11, 14, 7, 12, 6\}.
\]
%
Its elements pair up as $(3,14)$, $(10,7)$, $(5,12)$, $(11,6)$, each
pair summing to $0 \bmod 17$ --- this is $-\,g\omega^{j} = g\omega^{j+4}$
in action. Squaring gives $D^{(2)} = \{9, 15, 8, 2\}$, of size $4$,
which is exactly the coset
$\omega\cdot\langle\omega^2\rangle = 9\cdot\langle 13\rangle$ built from
the updated parameters.
\end{example}


Folding is the algebraic core of FRI: it maps a polynomial of degree
less than $d$, evaluated on $D$, to a polynomial of degree less than
$d/2$, evaluateend{example}


Folding is the algebraic core of FRI: it maps a polynomial of degree
less than $d$, evaluated on $D$, to a polynomial of degree less than
$d/2$, evaluated on $D^{(2)}$, using a single random challenge.
 
 
\subsection{Even-Odd Decomposition}
\label{subsec:even_odd}
 
\begin{definition}[Even-Odd Decomposition]
\label{def:even_odd}
Let $f(X) = \sum_{i=0}^{d-1} c_i X^i \in \Fp[X]$. Define
%
\[
f_{e}(Y) = \sum_{i \text{ even}} c_i \, Y^{i/2},
\qquad
f_{o}(Y) = \sum_{i \text{ odd}} c_i \, Y^{(i-1)/2}.
\]
%
Then $f$ decomposes uniquely as
%
\begin{equation}
\label{eq:even_odd}
f(X) = f_{e}(X^2) + X \cdot f_{o}(X^2),
\end{equation}
%
and both $f_e$ and $f_o$ have degree less than $\lceil d/2 \rceil$.
\end{definition}
 
The decomposition simply separates the coefficients of $f$ by the
parity of their index; the degree bound follows because each of the two
polynomials retains at most half of the $d$ coefficients. This is the
same identity that underlies the recursive step of the NTT
(\cref{subsec:ntt}), and it is the split used by FRI in the commit
phase~\cite[Section~3.9.1]{ethSTARK}.
 
 
\subsection{Recovering the Components from a Pair}
\label{subsec:recovering}
 
The prover holds $f$ in evaluation form, not in coefficient form, so
the decomposition must be computable from evaluations alone. Evaluating
\cref{eq:even_odd} at $x$ and at $-x$ gives the linear
system~\cite[Section~3.9.2]{ethSTARK}
%
\begin{equation}
\label{eq:fold_system}
\begin{aligned}
f(x)  &= f_{e}(x^2) + x \cdot f_{o}(x^2), \\
f(-x) &= f_{e}(x^2) - x \cdot f_{o}(x^2),
\end{aligned}
\end{equation}
%
in the two unknowns $f_{e}(x^2)$ and $f_{o}(x^2)$. Since $x \neq 0$ for
every $x \in D$, the system is nonsingular and solves to
%
\begin{equation}
\label{eq:fold_solve}
f_{e}(x^2) = \frac{f(x) + f(-x)}{2},
\qquad
f_{o}(x^2) = \frac{f(x) - f(-x)}{2x}.
\end{equation}
%
Both components at the point $x^2 \in D^{(2)}$ are thus determined by
the single pair $\bigl(f(x), f(-x)\bigr)$, which by
\cref{prop:domain}(2) is the pair of codeword entries at indices $i$
and $i + N/2$.
 
 
\subsection{The Fold}
\label{subsec:the_fold}
 
\begin{definition}[FRI Fold]
\label{def:fold}
Let $f$ have degree less than $d$ and let $\alpha$ be a challenge
sampled from a field $\mathbb{K} \supseteq \Fp$. The \emph{fold} of $f$
with respect to $\alpha$ is
%
\begin{equation}
\label{eq:fold}
f^{(\alpha)}(Y) := f_{e}(Y) + \alpha \cdot f_{o}(Y) ,
\end{equation}
%
a polynomial over $\mathbb{K}$ of degree less than $\lceil d/2 \rceil$.
\end{definition}
 
Combining \cref{eq:fold} with \cref{eq:fold_solve}, the codeword of
$f^{(\alpha)}$ on $D^{(2)}$ is computed from the codeword of $f$ on $D$
without any interpolation:
%
\begin{equation}
\label{eq:fold_evaluations}
f^{(\alpha)}(x^2)
  = \frac{f(x) + f(-x)}{2}
  + \alpha \cdot \frac{f(x) - f(-x)}{2x} .
\end{equation}
%
This is precisely the operation performed by our implementation: for
each $i < N/2$ it reads the entries at indices $i$ and $i + N/2$,
applies \cref{eq:fold_evaluations} with $x = g\omega^{i}$, and emits a
codeword of length $N/2$ over the domain $D^{(2)}$ described by the
updated parameters of \cref{eq:domain_update}.
 
A single round therefore halves two quantities at once: the degree
bound, by \cref{def:fold}, and the domain size, by
\cref{prop:domain}(3). The rate $\rho$ is invariant, and the folded
codeword is again a Reed-Solomon codeword of the same rate.
 
% alpha is sampled from F_{p^4} by the challenger; from round 0 onwards
% the codeword entries are extension field elements.
\begin{remark}[Arithmetic in the extension field]
\label{rem:ext_field_fold}
The initial codeword lies in $\Fp$, but the challenge $\alpha$ is
sampled from the quartic extension $\Fpk$
(\cref{subsec:field_ext}). By \cref{eq:fold}, the folded polynomial has
coefficients in $\Fpk$, and every subsequent layer is an
$\Fpk$-codeword. In our implementation this transition happens at the
first fold: the entries of the initial codeword are base field
elements, while all later layers hold extension field elements, and the
Merkle leaves grow accordingly from four to sixteen bytes.
\end{remark}
 
\begin{example}
\label{ex:fold_f17}
Continuing \cref{ex:domain_f17}, let
$f(X) = 1 + 2X + 3X^2 + 4X^3$ over $\F_{17}$, so $d = 4$, $N = 8$ and
$\rho = 1/2$. The codeword of $f$ on $D = \{3, 10, 5, 11, 14, 7, 12,
6\}$ is
%
\[
\bigl(f(x)\bigr)_{x \in D} = (6, 3, 8, 15, 16, 4, 8, 16).
\]
%
The decomposition of \cref{def:even_odd} gives $f_e(Y) = 1 + 3Y$ and
$f_o(Y) = 2 + 4Y$. With the challenge $\alpha = 5$,
%
\[
f^{(5)}(Y) = (1 + 3Y) + 5\,(2 + 4Y) = 11 + 6Y ,
\]
%
of degree $1 < 2$, as predicted by \cref{def:fold}. Applying
\cref{eq:fold_evaluations} to the four pairs
$\bigl(f(x), f(-x)\bigr)$ produces the codeword
$(14, 16, 8, 6)$ on $D^{(2)} = \{9, 15, 8, 2\}$, which is exactly the
evaluation of $11 + 6Y$ on that domain. Both the degree bound and the
domain size have halved, while the rate remains $2/4 = 1/2$.
\end{example}
 d on $D^{(2)}$, using a single random challenge.
 
 
\subsection{Even-Odd Decomposition}
\label{subsec:even_odd}
 
\begin{definition}[Even-Odd Decomposition]
\label{def:even_odd}
Let $f(X) = \sum_{i=0}^{d-1} c_i X^i \in \Fp[X]$. Define
%
\[
f_{e}(Y) = \sum_{i \text{ even}} c_i \, Y^{i/2},
\qquad
f_{o}(Y) = \sum_{i \text{ odd}} c_i \, Y^{(i-1)/2}.
\]
%
Then $f$ decomposes uniquely as
%
\begin{equation}
\label{eq:even_odd}
f(X) = f_{e}(X^2) + X \cdot f_{o}(X^2),
\end{equation}
%
and both $f_e$ and $f_o$ have degree less than $\lceil d/2 \rceil$.
\end{definition}
 
The decomposition simply separates the coefficients of $f$ by the
parity of their index; the degree bound follows because each of the two
polynomials retains at most half of the $d$ coefficients. This is the
same identity that underlies the recursive step of the NTT
(\cref{subsec:ntt}), and it is the split used by FRI in the commit
phase~\cite[Section~3.9.1]{ethSTARK}.
 
 
\subsection{Recovering the Components from a Pair}
\label{subsec:recovering}
 
The prover holds $f$ in evaluation form, not in coefficient form, so
the decomposition must be computable from evaluations alone. Evaluating
\cref{eq:even_odd} at $x$ and at $-x$ gives the linear
system~\cite[Section~3.9.2]{ethSTARK}
%
\begin{equation}
\label{eq:fold_system}
\begin{aligned}
f(x)  &= f_{e}(x^2) + x \cdot f_{o}(x^2), \\
f(-x) &= f_{e}(x^2) - x \cdot f_{o}(x^2),
\end{aligned}
\end{equation}
%
in the two unknowns $f_{e}(x^2)$ and $f_{o}(x^2)$. Since $x \neq 0$ for
every $x \in D$, the system is nonsingular and solves to
%
\begin{equation}
\label{eq:fold_solve}
f_{e}(x^2) = \frac{f(x) + f(-x)}{2},
\qquad
f_{o}(x^2) = \frac{f(x) - f(-x)}{2x}.
\end{equation}
%
Both components at the point $x^2 \in D^{(2)}$ are thus determined by
the single pair $\bigl(f(x), f(-x)\bigr)$, which by
\cref{prop:domain}(2) is the pair of codeword entries at indices $i$
and $i + N/2$.
 
 
\subsection{The Fold}
\label{subsec:the_fold}
 
\begin{definition}[FRI Fold]
\label{def:fold}
Let $f$ have degree less than $d$ and let $\alpha$ be a challenge
sampled from a field $\mathbb{K} \supseteq \Fp$. The \emph{fold} of $f$
with respect to $\alpha$ is
%
\begin{equation}
\label{eq:fold}
f^{(\alpha)}(Y) := f_{e}(Y) + \alpha \cdot f_{o}(Y) ,
\end{equation}
%
a polynomial over $\mathbb{K}$ of degree less than $\lceil d/2 \rceil$.
\end{definition}
 
Combining \cref{eq:fold} with \cref{eq:fold_solve}, the codeword of
$f^{(\alpha)}$ on $D^{(2)}$ is computed from the codeword of $f$ on $D$
without any interpolation:
%
\begin{equation}
\label{eq:fold_evaluations}
f^{(\alpha)}(x^2)
  = \frac{f(x) + f(-x)}{2}
  + \alpha \cdot \frac{f(x) - f(-x)}{2x} .
\end{equation}
%
This is precisely the operation performed by our implementation: for
each $i < N/2$ it reads the entries at indices $i$ and $i + N/2$,
applies \cref{eq:fold_evaluations} with $x = g\omega^{i}$, and emits a
codeword of length $N/2$ over the domain $D^{(2)}$ described by the
updated parameters of \cref{eq:domain_update}.
 
A single round therefore halves two quantities at once: the degree
bound, by \cref{def:fold}, and the domain size, by
\cref{prop:domain}(3). The rate $\rho$ is invariant, and the folded
codeword is again a Reed-Solomon codeword of the same rate.
 
% alpha is sampled from F_{p^4} by the challenger; from round 0 onwards
% the codeword entries are extension field elements.
\begin{remark}[Arithmetic in the extension field]
\label{rem:ext_field_fold}
The initial codeword lies in $\Fp$, but the challenge $\alpha$ is
sampled from the quartic extension $\Fpk$
(\cref{subsec:field_ext}). By \cref{eq:fold}, the folded polynomial has
coefficients in $\Fpk$, and every subsequent layer is an
$\Fpk$-codeword. In our implementation this transition happens at the
first fold: the entries of the initial codeword are base field
elements, while all later layers hold extension field elements, and the
Merkle leaves grow accordingly from four to sixteen bytes.
\end{remark}
 
\begin{example}
\label{ex:fold_f17}
Continuing \cref{ex:domain_f17}, let
$f(X) = 1 + 2X + 3X^2 + 4X^3$ over $\F_{17}$, so $d = 4$, $N = 8$ and
$\rho = 1/2$. The codeword of $f$ on $D = \{3, 10, 5, 11, 14, 7, 12,
6\}$ is
%
\[
\bigl(f(x)\bigr)_{x \in D} = (6, 3, 8, 15, 16, 4, 8, 16).
\]
%
The decomposition of \cref{def:even_odd} gives $f_e(Y) = 1 + 3Y$ and
$f_o(Y) = 2 + 4Y$. With the challenge $\alpha = 5$,
%
\[
f^{(5)}(Y) = (1 + 3Y) + 5\,(2 + 4Y) = 11 + 6Y ,
\]
%
of degree $1 < 2$, as predicted by \cref{def:fold}. Applying
\cref{eq:fold_evaluations} to the four pairs
$\bigl(f(x), f(-x)\bigr)$ produces the codeword
$(14, 16, 8, 6)$ on $D^{(2)} = \{9, 15, 8, 2\}$, which is exactly the
evaluation of $11 + 6Y$ on that domain. Both the degree bound and the
domain size have halved, while the rate remains $2/4 = 1/2$.
\end{example}
 